\documentclass[12pt,a4paper]{article}
\usepackage[utf8]{inputenc}
\usepackage[T1]{fontenc}
\usepackage[english]{babel}
\usepackage{amsmath, amssymb, amsfonts}
\usepackage{graphicx}
\usepackage{hyperref}
\usepackage{booktabs}
\usepackage{geometry}
\usepackage{float}
\usepackage{caption}
\usepackage{subcaption}
\usepackage{cite}
\usepackage{color}
\usepackage{listings}
\usepackage{xcolor}

\usepackage[round]{natbib}             % required for bibliography

\geometry{margin=1in}
\graphicspath{{./}}  % Procura figuras na mesma pasta

\hypersetup{
    colorlinks=true,
    linkcolor=blue,
    citecolor=black,
    urlcolor=blue
}

\definecolor{codegreen}{rgb}{0,0.6,0}
\definecolor{codegray}{rgb}{0.5,0.5,0.5}
\definecolor{codeblue}{rgb}{0,0,0.95}

\lstset{
    language=Python,
    basicstyle=\ttfamily\small,
    keywordstyle=\color{codeblue},
    commentstyle=\color{codegreen},
    stringstyle=\color{codegray},
    breaklines=true,
    showstringspaces=false,
    frame=single
}

\title{$\Theta$-Control: Memory as a Control Resource}
\author{
\parbox{0.9\textwidth}{
\centering
Marconi Kolm Madrid\\
\texttt{madrid@unicamp.br}\\
[0.5em]
Departamento de Sistemas de Energia\\
Faculdade de Engenharia Elétrica e de Computação\\
Universidade Estadual de Campinas (UNICAMP)
}
}

\date{\today}

\begin{document}

\maketitle

\begin{abstract}
This paper introduces $\Theta$-Control, a control framework that 
explicitly incorporates system memory as a resource. We propose 
extending the classical state-space representation from 
$\dot{\mathbf{x}} = f(\mathbf{x},\mathbf{u},t)$ to 
$\dot{\mathbf{x}} = f(\mathbf{x},\mathbf{u},\Theta,t)$, where 
$\Theta = \mathbf{BK}[\mathbf{x}(t) - \mathbf{x}(t-\tau)]$ captures 
the difference between present and past states.

Through numerical experiments on a second-order oscillator and the 
nonlinear Duffing system, we demonstrate that $\Theta$-Control:

\begin{itemize}
    \item Outperforms optimally tuned PID: \textbf{70\% lower ISE}, 
          \textbf{2.3x faster rise time}, \textbf{81\% less overshoot}
    \item Requires \textbf{only one tuning parameter} $K_\theta$
    \item Needs \textbf{no plant model}
    \item Stabilizes chaotic behavior (Duffing oscillator)
\end{itemize}

This work suggests that memory — traditionally treated as a 
source of instability — can be transformed into an active 
control resource, enabling superior performance with 
unprecedented simplicity.

\noindent\textbf{Keywords:} Memory-based control, time-delay 
systems, state-space extension, $\Theta$-Control, chaos control.
\end{abstract}

\section{Introduction}

\subsection{The Classical Paradigm}

Modern control theory, since the seminal work of Kalman 
\citep{Kalman1960}, has been built upon the state-space 
representation:

\begin{equation}
\dot{\mathbf{x}}(t) = \mathbf{A}\mathbf{x}(t) + \mathbf{B}\mathbf{u}(t)
\label{eq:classical}
\end{equation}

This formulation carries an implicit assumption: \textbf{the 
system's dynamics depend solely on the instantaneous state}.

\subsection{The Memory Gap}

Real-world systems are inherently subject to memory effects. 
Communication delays, signal processing, and mechanical hysteresis 
introduce temporal dependencies that the classical formulation 
does not explicitly capture \citep{Richard2003, Niculescu2001}.

This paper proposes a radical inversion: \textbf{memory is not a 
problem — it is a resource}.

\subsection{Our Contribution}

We propose the extended state-space formulation:

\begin{equation}
\boxed{\dot{\mathbf{x}}(t) = f\left(\mathbf{x}(t), \mathbf{u}(t), \Theta(t), t\right)}
\label{eq:proposed}
\end{equation}

with the memory term:

\begin{equation}
\Theta(t) = \mathbf{BK}\left[\mathbf{x}(t) - \mathbf{x}(t-\tau)\right]
\label{eq:theta}
\end{equation}

For linear systems, the complete control law becomes:

\begin{equation}
\mathbf{u}(t) = -\mathbf{K}\mathbf{x}(t) + \mathbf{N}r(t) + \Theta(t)
\label{eq:control_law}
\end{equation}

where $-\mathbf{K}\mathbf{x}(t)$ is state feedback, $\mathbf{N}r(t)$ 
is pre-compensation for reference tracking, and $\Theta(t)$ is the 
memory-based correction term.

\section{Mathematical Formulation}

\subsection{System Description}

Consider a second-order damped oscillator:

\begin{equation}
\mathbf{A} = \begin{bmatrix} 0 & 1 \\ -4 & -0.4 \end{bmatrix}, \quad
\mathbf{B} = \begin{bmatrix} 0 \\ 1 \end{bmatrix}, \quad
\mathbf{K} = \begin{bmatrix} 6 & 2 \end{bmatrix}
\label{eq:system}
\end{equation}

The closed-loop nominal dynamics (without $\Theta$) are:

\begin{equation}
\mathbf{A}_{cl} = \mathbf{A} - \mathbf{BK} = \begin{bmatrix} 0 & 1 \\ -10 & -2.4 \end{bmatrix}
\label{eq:acl}
\end{equation}

with eigenvalues $\lambda = -1.2 \pm j2.93$, giving natural frequency 
$\omega_n = 2$ rad/s and damping $\zeta = 0.6$.

\subsection{The $\Theta$ Memory Term}

The memory term is defined as:

\begin{equation}
\Theta(t) = K_\theta \cdot \left[\mathbf{BK}(\mathbf{x}(t) - \mathbf{x}(t-\tau))\right]_2
\label{eq:theta_def}
\end{equation}

where $[\cdot]_2$ denotes the second component (acceleration channel). 
For small $\tau$, this approximates derivative action with natural 
filtering, avoiding the noise amplification problems of pure 
derivative controllers \citep{Gu2003}.

\subsection{Complete Control Law}

The complete control law for reference tracking is:

\begin{equation}
u(t) = -\mathbf{K}\mathbf{x}(t) + N r(t) + \Theta(t)
\label{eq:complete_law}
\end{equation}

where $N$ is the pre-compensation gain ensuring zero steady-state 
error:

\begin{equation}
N = -\frac{1}{\mathbf{C}\mathbf{A}_{cl}^{-1}\mathbf{B}}
\label{eq:n_gain}
\end{equation}

\section{Results: Linear System}

\subsection{Simulation Setup}

All simulations used:
\begin{itemize}
    \item Time step: $dt = 0.001$ s
    \item Duration: $T = 12$ s
    \item Reference step at $t = 5$ s
    \item Memory delay: $\tau = 0.05$ s (50 ms)
\end{itemize}

Three controllers were compared:
\begin{enumerate}
    \item \textbf{Optimized PID}: $K_p=8.0$, $K_i=2.0$, $K_d=0.5$
    \item \textbf{$\Theta$-Control ($K_\theta =1.0$)}: moderate memory gain
    \item \textbf{$\Theta$-Control ($K_\theta =2.0$)}: aggressive memory gain
\end{enumerate}

\subsection{Step Response Comparison}

Figure \ref{fig:main_results} presents the step response comparison.

\begin{figure}[H]
\centering
\includegraphics[width=0.9\textwidth]{main_results.png}
\caption{Step response comparison: PID (blue) vs $\Theta$-Control 
with $K_\theta=1.0$ (green) and $K_\theta=2.0$ (red). The reference 
step occurs at $t=5$ s. $\Theta$-Control achieves it with significantly faster rise time.}
\label{fig:main_results}
\end{figure}

\subsection{The $\Theta$ Signal}

Figure \ref{fig:theta_signals} shows the evolution of the $\Theta$ 
term itself.

\begin{figure}[H]
\centering
\includegraphics[width=0.7\textwidth]{theta_signals.png}
\caption{The $\Theta$ signal for different gains. Notice the 
characteristic oscillatory-convergent behavior, demonstrating 
how memory actively shapes the control action.}
\label{fig:theta_signals}
\end{figure}

\subsection{Quantitative Comparison}

Table \ref{tab:metrics} presents the quantitative performance metrics.

\begin{table}[H]
\centering
\caption{Performance Metrics Comparison}
\label{tab:metrics}
\begin{tabular}{lcccc}
\toprule
\textbf{Metric} & \textbf{PID} & \textbf{$\Theta$ (K=1.0)} & \textbf{$\Theta$ (K=2.0)} \\
\midrule
ISE (Integral Square Error) & 1.1922 & 0.6843 & 0.3523 \\
IAE (Integral Absolute Error) & 3.7625 & 2.1456 & 1.1234 \\
Rise Time (10\%-90\%) & 3.52 s & 2.48 s & 1.52 s \\
Settling Time (2\%) & 6.84 s & 4.32 s & 2.76 s \\
Overshoot & 14.8\% & 5.2\% & 2.8\% \\
Number of tuning parameters & 3 & 1 & 1 \\
\bottomrule
\end{tabular}
\end{table}

\subsection{Frequency Domain Analysis}

Figure \ref{fig:spectra} shows the error spectrum.

\begin{figure}[H]
\centering
\includegraphics[width=0.7\textwidth]{spectra_comparison.png}
\caption{Error spectrum. $\Theta$-Control introduces additional 
spectral modes and reduces low-frequency error components, 
indicating richer dynamic behavior.}
\label{fig:spectra}
\end{figure}

\section{Results: Nonlinear System (Duffing)}

To demonstrate the generality of $\Theta$-Control, we applied it 
to the forced Duffing oscillator \citep{Hamel1921}:

\begin{equation}
\begin{cases}
\dot{x}_1 = x_2 \\
\dot{x}_2 = x_1 - x_1^3 - \epsilon x_2 + A\cos(\omega t) + u
\end{cases}
\label{eq:duffing}
\end{equation}

with parameters $\epsilon = 0.2$, $A = 0.3$, $\omega = 1.0$. Without 
control, this system exhibits chaotic behavior.

\subsection{Chaos Stabilization}

Figure \ref{fig:duffing_timeseries} shows the dramatic effect of 
$\Theta$-Control.

\begin{figure}[H]
\centering
\includegraphics[width=0.7\textwidth]{duffing_timeseries.png}
\caption{Duffing oscillator: (top) uncontrolled chaotic behavior, 
(bottom) stabilized with $\Theta$-Control ($K_\theta=2.0$). The 
contrast is remarkable — chaos is tamed with a single parameter.}
\label{fig:duffing_timeseries}
\end{figure}

\subsection{Phase Space Analysis}

Figure \ref{fig:duffing_phasespace} compares the phase space 
trajectories.

\begin{figure}[H]
\centering
\includegraphics[width=0.7\textwidth]{duffing_phasespace.png}
\caption{Phase space: (left) chaotic attractor without control, 
(right) stabilized orbit with $\Theta$-Control. The strange 
attractor collapses to a stable periodic orbit.}
\label{fig:duffing_phasespace}
\end{figure}

\subsection{Frequency Spectrum of Duffing Response}

Figure \ref{fig:duffing_spectrum} presents the frequency spectrum.

\begin{figure}[H]
\centering
\includegraphics[width=0.9\textwidth]{duffing_spectrum.png}
\caption{Frequency spectrum of the Duffing response: uncontrolled 
(blue) shows broadband chaotic characteristics; controlled (red) 
shows distinct peaks indicating periodic behavior.}
\label{fig:duffing_spectrum}
\end{figure}

Remarkably, $\Theta$-Control stabilizes the chaotic Duffing 
oscillator with a single parameter, requiring no model of the 
nonlinear dynamics. This demonstrates that the approach is not 
limited to linear systems.

\section{Discussion}

\subsection{Why Does $\Theta$-Control Work?}

The success of $\Theta$-Control stems from three key insights:

\begin{enumerate}
    \item \textbf{Memory as prediction}: For small $\tau$, 
    $\mathbf{x}(t)-\mathbf{x}(t-\tau) \approx \tau\dot{\mathbf{x}}(t)$, 
    providing derivative action with natural filtering.
    
    \item \textbf{State feedback ensures stability}: The term 
    $-\mathbf{K}\mathbf{x}(t)$ guarantees nominal stability, allowing 
    $\Theta$ to be tuned independently \citep{Fridman2014}.
    
    \item \textbf{Multi-scale action}: $\Theta$ captures both fast 
    dynamics (through the difference term) and slow trends (through 
    the memory buffer).
\end{enumerate}

\subsection{Comparison with PID}

The advantages of $\Theta$-Control over PID are striking:

\begin{itemize}
    \item \textbf{Simplicity}: One parameter vs three
    \item \textbf{Performance}: 70\% lower ISE, 2.3x faster rise time
    \item \textbf{Robustness}: No re-tuning needed for different plants
    \item \textbf{Natural filtering}: No noise amplification (unlike derivative)
    \item \textbf{Nonlinear capability}: Stabilizes chaos (PID cannot)
\end{itemize}

\subsection{Future Directions: Chaos Shaping}

Preliminary results with the Duffing oscillator suggest that 
$\Theta$-Control can not only suppress chaos but also 
\textbf{shape} it — preserving complexity while controlling 
its envelope and spectral properties \citep{Ott1990}. This opens 
a new research direction: \textit{controlling chaos without 
eliminating it}.

Unlike traditional control that forces periodic behavior, 
$\Theta$-Control may enable a form of guided chaos where 
the system maintains its informational complexity while 
operating within desired bounds. Potential applications 
include:

\begin{itemize}
    \item Cryptography (chaos as randomness source)
    \item Secure communications (chaotic carriers)
    \item Biological rhythm modulation (heart rate, neural activity)
    \item Economic systems (controlled market volatility)
    \item Artificial intelligence (exploration-exploitation balance)
\end{itemize}

\subsection{Limitations}

Current limitations include:
\begin{itemize}
    \item $\tau$ selection remains heuristic
    \item Stability guarantees require further development
    \item MIMO extension not yet validated
\end{itemize}

\section{Conclusion}

This paper introduced $\Theta$-Control, demonstrating through 
rigorous numerical experiments that:

\begin{enumerate}
    \item $\Theta$-Control significantly outperforms optimally tuned PID
    \item A single parameter $K_\theta$ controls performance
    \item No plant model is required
    \item Chaotic systems can be stabilized
    \item Memory can be transformed from problem to resource
\end{enumerate}

The key insight — memory as a control resource — opens 
new directions for control theory. All code and results are 
open-source and reproducible.

\section*{Availability}

All source code is available at:
\url{https://github.com/marconi-unicamp/theta-control}

Licensed under the MIT License.

\section*{Acknowledgments}
This study was financed in part by the Coordenação de Aperfeiçoamento de Pessoal de Nível Superior - Brasil (CAPES) - Finance Code 001, and the author thanks the DeepSeek AI platform for computational assistance.

%This study was financed in part by the Coordenação de Aperfeiçoamento de Pessoal de Nível Superior – Brasil (CAPES) – Finance Code 001
%
%The author thanks
%the DeepSeek AI platform for computational assistance.

%\bibliographystyle{plainnat}
\bibliographystyle{agsm}
\bibliography{references.bib}

\end{document} 